\chapter{Couche 6 — Le durcissement des conteneurs}

Le locataire est \cmd{root} dans sa machine : c'est le service vendu. Toute la
question est de faire en sorte que ce \cmd{root}-là ne soit pas \cmd{root} sur
le nœud.

% -----------------------------------------------------------------------------
\mesure{Les capacités retirées}{mesure 18 · \fichier{docker\_ops.py}}

\menace
Par défaut, Docker accorde à un conteneur une quinzaine de \emph{capacités} du
noyau Linux. Certaines sont anodines ; d'autres permettent de modifier le
réseau de l'hôte, de monter des systèmes de fichiers ou de charger des modules.

\parade
On retire \textbf{tout}, puis on rend uniquement ce dont une machine Linux
ordinaire a besoin pour démarrer ses services. C'est une liste blanche, pas une
liste noire.

\danslecode
\begin{lstlisting}[language=bash]
--cap-drop ALL
--cap-add CHOWN --cap-add DAC_OVERRIDE --cap-add FOWNER --cap-add FSETID
--cap-add SETGID --cap-add SETUID --cap-add SYS_CHROOT
--cap-add NET_BIND_SERVICE --cap-add KILL --cap-add AUDIT_WRITE
\end{lstlisting}

\begin{center}
\small
\begin{tabularx}{\textwidth}{@{}p{3.3cm}X@{}}
\toprule
\textbf{Capacité gardée} & \textbf{Pourquoi elle est nécessaire} \\
\midrule
\cmd{CHOWN}, \cmd{FOWNER}, \cmd{FSETID} & Créer les répertoires de service avec
les bons propriétaires et droits (\cmd{/root/.ssh}, \cmd{/run/sshd}). \\
\cmd{DAC\_OVERRIDE} & Permettre à \cmd{root} d'écrire là où les droits ne le
prévoient pas --- comportement normal de \cmd{root}. \\
\cmd{SETGID}, \cmd{SETUID} & \cmd{login} et \cmd{sshd} changent d'identité pour
ouvrir une session : sans elles, aucune connexion. \\
\cmd{SYS\_CHROOT} & \cmd{sshd} isole son processus non privilégié par
\cmd{chroot} pendant l'authentification. \\
\cmd{NET\_BIND\_SERVICE} & Écouter sur un port inférieur à 1024 --- le 22. \\
\cmd{KILL} & \cmd{supervisord} doit pouvoir arrêter et relancer ses services. \\
\cmd{AUDIT\_WRITE} & PAM écrit un enregistrement d'audit à chaque connexion. \\
\bottomrule
\end{tabularx}
\end{center}

\begin{notion}{Ce qui est retiré, et qui compte}
\cmd{SYS\_ADMIN} (monter des systèmes de fichiers --- la capacité la plus
dangereuse, souvent qualifiée de « nouveau \cmd{root} »), \cmd{NET\_ADMIN}
(reconfigurer le réseau, poser une interface en mode espion),
\cmd{SYS\_MODULE} (charger un module noyau), \cmd{SYS\_PTRACE} (inspecter la
mémoire d'autres processus), \cmd{MKNOD} (créer un fichier de périphérique,
donc accéder au disque brut).

Chacune de ces capacités est une voie d'évasion documentée. Les retirer ne
retire rien à l'usage légitime d'une machine louée.
\end{notion}

\limite
Les capacités cloisonnent le noyau, elles ne le remplacent pas. Une
vulnérabilité du noyau lui-même traverse la frontière.

% -----------------------------------------------------------------------------
\mesure{Pas d'élévation de privilèges}{mesure 19 · \fichier{docker\_ops.py}}

\menace
Un binaire \cmd{setuid} à l'intérieur du conteneur peut rendre des privilèges
qu'on vient justement de retirer. Beaucoup d'images en contiennent sans le
savoir (\cmd{ping}, \cmd{mount}, \cmd{su}).

\parade
\cmd{-{}-security-opt no-new-privileges:true} : le noyau refuse qu'un processus
du conteneur gagne des privilèges qu'il n'avait pas au départ, quel que soit le
bit \cmd{setuid} du binaire exécuté.

\danslecode
Dans la constante \cmd{CONTAINER\_HARDENING}, appliquée à \emph{toutes} les
machines louées sans exception --- ce n'est pas une option du formulaire.

\limite
Le drapeau empêche de \emph{gagner} des privilèges. Il ne retire pas ceux que
\cmd{root} possède déjà dans le conteneur.

% -----------------------------------------------------------------------------
\mesure{Les limites de ressources}{mesure 20 · \fichier{docker\_ops.py}}

\menace
Un locataire lance une bombe à fourche (\cmd{:(){ :|:\& };:}), un mineur de
cryptomonnaie ou une compilation massive. Sans limite, il prend tout le nœud et
toutes les autres machines deviennent inutilisables. C'est le déni de service
le plus facile à monter : il suffit d'être client.

\parade
Trois plafonds, posés au lancement.

\danslecode
\begin{lstlisting}[language=bash]
--pids-limit 512            # nombre maximal de processus : arrete la bombe a fourche
--cpus 1                    # un coeur au plus
--memory 256m | 1g          # selon le mode, terminal ou bureau
--memory-reservation ...    # le minimum garanti en cas de contention
\end{lstlisting}

La distinction entre \cmd{-{}-memory} et \cmd{-{}-memory-reservation} est celle
que le cours impose : la \emph{limite} est un plafond dur que le conteneur ne
peut pas dépasser, la \emph{réservation} est le minimum qui lui reste garanti
quand le nœud est sous pression. Sans réservation, une machine bien élevée peut
être étranglée par une voisine gourmande, même si celle-ci respecte sa limite.

\begin{retenir}
\cmd{-{}-pids-limit} est la plus sous-estimée des trois : une bombe à fourche ne
consomme ni beaucoup de CPU ni beaucoup de mémoire par processus, mais épuise
la table des processus du \emph{noyau}, qui est partagée avec l'hôte. Les
limites de CPU et de mémoire ne l'arrêtent pas.
\end{retenir}

\limite
Les limites valent par conteneur. Trois machines par utilisateur, c'est trois
fois la limite --- d'où le quota.

% -----------------------------------------------------------------------------
\mesure{Un réseau dédié}{mesure 21 · \fichier{config.py}, \fichier{docker\_ops.py}}

\menace
Sur le réseau \cmd{bridge} par défaut de Docker, tous les conteneurs se voient.
Un locataire peut donc scanner et atteindre les machines des autres locataires
--- et les services de la plateforme elle-même.

\parade
Les machines louées sont placées sur un réseau Docker \textbf{propre au
déploiement}, créé à la demande.

\danslecode
\begin{lstlisting}[language=Python]
CONTAINER_NETWORK = os.environ.get("INSACLOUD_NETWORK", f"insacloud_{CLUSTER}_net")
\end{lstlisting}

\begin{piege}{\cmd{-{}-internal} aurait coupé Internet aux locataires}
La première version créait le réseau avec \cmd{-{}-internal}, qui l'isole
totalement --- y compris de l'extérieur. L'intention était bonne, l'effet
désastreux : plus aucune machine louée ne pouvait faire \cmd{apt-get install},
ce qui est précisément ce pour quoi on loue une machine.

Le défaut a été vu avant livraison, en essayant réellement d'installer un
paquet dans une machine fraîchement louée. C'est le genre d'erreur qu'aucune
relecture ne détecte, parce que le code est exactement conforme à l'intention.
\end{piege}

\limite
Le réseau dédié sépare les machines louées des services de la plateforme ; il
ne sépare pas les locataires \emph{entre eux}. Un réseau par machine serait
plus strict, au prix d'une complexité de gestion sans commune mesure avec le
gain à cette échelle.

% -----------------------------------------------------------------------------
\mesure{L'identité par étiquette}{mesure 22 · \fichier{docker\_ops.py}}

\menace
Le Faucheur détruit les conteneurs qui ne correspondent à aucune location
active. S'il les reconnaît par leur \emph{nom}, une collision de préfixe lui
fait détruire autre chose --- par exemple les services de la plateforme, qui
s'appellent \cmd{insacloud-web}, \cmd{insacloud-db}, \cmd{insacloud-faucheur}.

\parade
Trois étiquettes Docker, apposées à chaque machine, qui servent d'identité.
Jamais le nom.

\danslecode
\begin{lstlisting}[language=Python]
MACHINE_LABEL = "insacloud.role"      # = rented-machine
CLUSTER_LABEL = "insacloud.cluster"   # = ce deploiement
WORKER_LABEL  = "insacloud.worker"    # = le noeud qui l'heberge

MACHINE_FILTERS = [f"label={MACHINE_LABEL}={MACHINE_LABEL_VALUE}",
                   f"label={CLUSTER_LABEL}={CLUSTER}"]
\end{lstlisting}

Les deux premiers filtres sont \emph{combinés} : une machine doit être une
machine louée \textbf{et} appartenir à ce déploiement. Sans le second, le
Faucheur d'une plateforme détruirait les machines d'une autre tournant sur le
même démon --- ce qui est arrivé en développement.

Le troisième, \cmd{insacloud.worker}, a été ajouté plus tard : il dit sur quel
nœud la machine a été créée. Sur un parc de machines virtuelles, la frontière
des démons suffit ; dès que deux nœuds partagent un démon, elle disparaît, et
sans ce label la réconciliation de l'un détruisait les machines de l'autre.

\begin{retenir}
Ce défaut n'a pas été trouvé en relisant le code --- il a été trouvé en faisant
tourner un parc à deux nœuds. Une étiquette est une \emph{donnée portée par
l'objet lui-même} ; un nom est une convention que rien ne garantit. Pour une
opération destructive, la différence n'est pas théorique.
\end{retenir}

\limite
Un attaquant ayant accès au démon Docker peut poser lui-même ces étiquettes sur
un conteneur à lui, et le faire adopter par la plateforme. Mais, une fois de
plus : qui a le démon a déjà tout.

% -----------------------------------------------------------------------------
\mesure{La commande Docker, sans interpréteur}{mesure 23 · \fichier{docker\_ops.py}}

\menace
L'énoncé impose de lancer les machines par la commande \cmd{docker run}. Écrite
en concaténant des chaînes et passée à un shell, une telle commande est une
injection en puissance : une distribution nommée
\cmd{alpine; rm -rf /} suffirait.

\parade
La commande est construite en \textbf{liste d'arguments} et exécutée sans
shell. Aucun interpréteur ne relit jamais les valeurs.

\danslecode
\begin{lstlisting}[language=Python]
res = subprocess.run(cmd, capture_output=True, text=True, timeout=timeout)
# pas de shell=True : les arguments sont passes en liste, donc jamais
# reinterpretes par un interpreteur de commandes.
\end{lstlisting}

Et en amont, les seules valeurs d'origine utilisateur --- durée, distribution,
mode --- sont validées : la durée est un entier borné, la distribution et le
mode sont choisis dans des dictionnaires fermés. Une valeur hors de ces
ensembles n'atteint jamais \cmd{docker}.

\begin{notion}{Deux barrières plutôt qu'une}
La validation en amont suffirait. L'absence de shell suffirait. On fait les
deux, parce que la première dépend d'un développeur qui n'oublie pas un
chemin, et la seconde d'une propriété du système. Quand l'une cède, l'autre
tient.
\end{notion}

\limite
L'exécution sans shell ne protège pas d'un argument \emph{valide} mais
malveillant --- par exemple un chemin de montage arbitraire, si un jour on
laissait l'utilisateur en choisir un. La liste des arguments est entièrement
construite par la plateforme : c'est ce qui doit le rester.
